\label{ch:discussion}



\subsection{Feedback and Limitations}
\label{sec:feedback-and-limitations}
\textbf{Skrives om}

The feedback from both projects gave useful insight into the usability of 
MLguide and what could be improved. However, two limitations applied.

The first concerned the feedback format. Feedback was collected as part of
the students' project reports, which fit naturally within the course structure
and allowed students to reflect on MLguide in the context of their own problem
and data. Since each group worked on a different problem, an open format seemed
more suitable than a fixed set of questions. The trade-off was that responses
varied in detail and focus, making it more difficult to draw systematic conclusions.

The second limitation was that the feedback arrived quite late in the development
period. This was probably unavoidable, as a working system had to be in place 
before it could be tested. The feedback led to some adjustments, but larger 
changes such as new features or structural redesigns were not feasible within 
the remaining time. These can instead be left as directions for future work.